\subsection{Selecting SSL configurations from gradient logs}\label{sec:ssl}
Joint-embedding self-supervised learning (SSL) trains representations to agree across augmented views. Agreement alone can admit collapsed solutions that map different images to the same representation \citep{chen2021simsiam}; training loss is therefore an unreliable guide to representation quality \citep{garrido2023rankme}. We test whether gradient-norm MG can guide configuration choice without extracting embeddings from each evaluated model. The MNIST experiment compares 29 configurations across four objectives using a 784--256--256 MLP encoder. We choose the log and score direction from labeled development runs, then freeze the rule for five test seeds. Applying it requires no labels.

MG selects configurations with low average regret and ranks them more accurately than encoder RankMe, the tested $\alpha$-ReQ slope and LiDAR implementations \citep{garrido2023rankme,agrawal2022alphareq,thilak2024lidar}. Table~\ref{tab:ssl_campaign} reports correlations and selection regret. Embedding spread has the highest mean correlation; several simple scalar features remain competitive. The benefit is useful selection from a saved gradient log without an embedding pass. \appref{app:campaign} details the single-dataset scope, proxy implementations and uncertainty.

\subsection{Warning of loss spikes in transformer training}
In a two-layer transformer trained on modular addition, MG of log-loss warns of most test spikes within the specified horizon. It has higher AUROC than the tested attention-logit and attention-entropy monitors. Roughness detects more events at the same false-positive rate for less computation (Table~\ref{tab:spikes_campaign}). This result supports detecting the slingshot regime from scalar logs. Preventing spikes and transferring the rule to language models remain untested (\appref{app:campaign}).
